\begin{abstract}
A Bayesian calibration of a cone-crusher population-balance engine with a structural discrepancy ties a line in the
closed-side setting on ten industrial surveys, and ten surveys cannot say when its physics would pay. We derive its prediction
risk given the scalars exactly: least squares on the discrepancy regressors plus the engine's departure from their plane, so the
engine gains over a regression through a shrinkage of the plane, which fades as surveys accumulate, and through its curvature,
which persists but helps only when aligned with the plant's and less than twice as large. We give the risk-optimal discrepancy
scale, the crossing size against the line, how the Kennedy-O'Hagan, orthogonal and scaled priors relate, and a floor on the
capacity base's contraction at any number of surveys at a fixed speed. A virtual-plant study of
\cellsTotal{} cells of \repsPerCell{} replicates matches the risk (largest $|z|$ \zMax{} of \zNum) and the sign change of the gain
at half the engine's curvature. On a plant with the published power-law capacity, an engine of the plant's form keeps an error
below least squares at 160 surveys (relative RMSE \capRoXBig{} under extrapolation), bounded by the plant's curvature, \rOverSigma\pct{}
of the noise; the shipped affine capacity cannot. With proper priors the orthogonal prior raises the power RMSE from
\rVzeroSTenKwKo{} to \rVzeroSTenKwOrth~kW at ten surveys. A product-size measurement at 2\pct{} noise adds more information than
doubling the surveys, and ten surveys already give \ceilShare\pct{} of the possible narrowing of the capacity base. \sbcAbstract{}
The virtual plants are models, not crushers.
\end{abstract}

\begin{IEEEkeywords}
Cone crusher, Bayesian calibration, model discrepancy, prediction risk, ridge regression, identifiability, experimental design,
simulation-based calibration, comminution.
\end{IEEEkeywords}
